\section{Algorithm 2: on-the-fly liveness}
\label{sec:otf}

Algorithm~\ref{alg:gfp} has one structural defect, and it is not a constant factor. Phase~1 runs
to completion before Phase~3 begins, so the cost of deciding \texttt{p --> q} is a function of the
\emph{model} and never of the \emph{query}: a leads-to property violated three transitions from the
initial state still pays for every reachable symbolic state. This section removes that defect
without weakening anything proved in \S\ref{sec:soundness}.

\subsection{Why this looks harder than it is}

$EG$ is a \emph{greatest} fixpoint, so the textbook iteration starts at $\top$ and descends, and an
intermediate iterate satisfies $W \supseteq F(W)$ --- it is a \textbf{pre}-fixpoint, an
over-approximation, and therefore says nothing. ``$W_v \neq \emptyset$'' cannot settle a violation,
so there is no early exit; and the moment the graph underneath grows, the values stop descending.

But a violation never needed the fixpoint. It needs a \emph{certificate}, and by Knaster--Tarski
the certificates of a greatest fixpoint are exactly its \textbf{post}-fixpoints --- the same
instrument used once already, in the proof of Theorem~\ref{thm:sound-holds}, here in the other role.

\begin{theorem}[A pass over a fragment is a certificate]
\label{thm:fragment}
Let $G$ be any subgraph explored so far and $S$ any set of \textbf{expanded} $\lnot q$ vertices of
$G$. Define $F_S$ exactly as $F$ but with $W_u$ pinned to $\emptyset$ for $u \notin S$. Then
\[
   \nu F_S \;\subseteq\; \nu F \;=\; W^\infty,
\]
where $F$ is the operator on the \emph{finished} graph. Consequently $\nu F_S(v) \neq \emptyset$ at
some $v \models p \wedge \lnot q$ settles the query as \textsc{violated}, with exactly the verdict
Algorithm~\ref{alg:gfp} would have produced. \proved{}
\end{theorem}

\begin{proof}
$F_S$ is obtained from $F$ by deleting disjuncts: every $\Pre_e(W_{v'})$ with $v' \notin S$ becomes
$\Pre_e(\emptyset) = \emptyset$, and no term is added. So $F_S(X) \subseteq F(X)$ pointwise. Both
are monotone, hence $\nu F_S \subseteq \nu F$: $\nu F_S$ is a post-fixpoint of $F_S$, hence of $F$,
hence below the greatest post-fixpoint of $F$.
\end{proof}

Taking $S$ to be every expanded $\lnot q$ vertex of the \emph{finished} graph gives $F_S = F$, so
the final pass \emph{is} Algorithm~\ref{alg:gfp}: nothing is given up when the property holds. No
value has to descend across passes, because each pass is a self-contained certificate search.

\begin{remark}[The one restriction]
\label{rem:expanded}
A vertex may enter $S$ only once \textbf{expanded}. The stopping term reads the \emph{absence} of
successors ($\mathrm{adj}(v) = \emptyset$), and an unexpanded vertex has an empty adjacency for the
wrong reason; admitting one would \emph{add} a disjunct rather than delete one. Every other term
survives a partial graph untouched, because the $\Pre$ disjuncts only ever need edges that
\textbf{exist}, never the absence of one. Absence of edges is dangerous, presence is not --- that
asymmetry is what makes backward work on a half-built graph legitimate at all.
\end{remark}

\begin{remark}[A capped pass proves nothing]
\label{rem:capped}
Theorem~\ref{thm:fragment} is about the \emph{limit}. An iterate stopped early --- by a
federation-width cap or an iteration bound --- is a pre-fixpoint and is neither an over- nor an
under-approximation. Reading such a value as ``no certificate found'' converts an exhausted budget
into a silent \textsc{holds}. The implementation reports \textsc{unusable} instead, and the
differential harness scores it as the absence of a verdict.
\end{remark}

\subsection{The algorithm}
\label{sec:otf:alg}

The forward search is a \textbf{candidate generator}; the backward operator is the \textbf{judge}.
That division is what lets the search run on the fully subsumed graph and still be exact. The
shipping forward engine cannot do this: it must \emph{refuse} subsumption edges whenever it argues
that a run continues forever, because a redirected edge closes cycles the concrete system cannot
execute (\S\ref{sec:liveness}). Algorithm~2 follows \textbf{every} edge and lets an exact confined
pass adjudicate.

\begin{algorithm}[h]
\caption{On-the-fly leads-to. Returns \textsc{holds} or \textsc{violated}.}
\label{alg:otf}
\begin{algorithmic}[1]
\State \textbf{Forward store.} $aLU$ subsumption, plus an \emph{exact} convex-union merge: a
       newcomer is absorbed into a stored zone whenever their union is exactly one zone
       (\S\ref{sec:otf:merge}).
\State \textbf{Frontier.} Breadth-first, except that a newly generated $p \wedge \lnot q$ state
       jumps the queue --- until \textsc{AccBudget} of them have been refuted, after which the
       order reverts to breadth-first (\S\ref{sec:otf:order}).
\State \textbf{At generation.} When a state $v$ is interned, if $v \models p \wedge \lnot q$ and
       $\mathrm{timeDiverges}(v)$ or $\mathrm{hasDeadlock}(v)$, report \textsc{violated}
       immediately (\S\ref{sec:otf:local}).
\State \textbf{At an accepting vertex} $u$ popped from the frontier: expand $u$; if
       $\mathrm{adj}(u) = \emptyset$ report \textsc{violated}; otherwise run a $\lnot q$ DFS from
       $u$ following \emph{every} edge. On reaching a stopping vertex or a vertex on the DFS
       stack, \textbf{confirm}: run the confined descending iteration of
       Theorem~\ref{thm:fragment} over the stack. Non-empty at $u$ $\Rightarrow$
       \textsc{violated}. Refuted $\Rightarrow$ continue, up to \textsc{RefuteBudget}
       confirmations.
\State \textbf{On exhaustion} of the frontier: one exact pass over the accepting closure
       (\S\ref{sec:otf:closure}). Empty everywhere $\Rightarrow$ \textsc{holds}.
\end{algorithmic}
\end{algorithm}

Steps 2, 3 and the budgets are pure \emph{ordering and early-exit} devices: they change which
witness is found first and how soon, never which states are reachable, so
Theorem~\ref{thm:fragment} and \S\ref{sec:soundness} carry over unchanged. Step 5 is what makes
the algorithm complete.

\subsection{The forward store: subsumption plus an exact merge}
\label{sec:otf:merge}

The graph is stored under $aLU$ subsumption, and additionally a newcomer is merged into a stored
zone whenever \texttt{dbm\_try\_union} succeeds --- that is, whenever the union of the two zones
\emph{really is a single zone}. The stored zone then denotes exactly the set the two it replaced
denoted: no valuation is added and none removed.

That exactness is what makes it usable here. An over-approximating merge (a convex hull, say) would
be sound only for concluding \textsc{holds}, because a certificate could then rest on valuations
the system never reaches. \texttt{try\_union} adds nothing, so it cannot invent a certificate and
the \textsc{violated} direction survives.

This matters more than it sounds. On \texttt{grc4} the $aLU$ antichain alone stores $18.32$ zones
per reachable discrete state, against $1.00$--$1.44$ on every other family; with the merge it
stores $6.90$, and the whole state space drops from $2345$ to $883$ symbolic states.
\measured{} The $14.5$--$18$ zones per key were never intrinsic to the model: they were an
antichain full of zones whose pairwise unions happen to be convex.

\subsection{Local witnesses, tested at generation}
\label{sec:otf:local}

Two of the three stopping disjuncts need \emph{no adjacency}: $\mathrm{timeDiverges}$ and
$\mathrm{hasDeadlock}$ are functions of a vertex's own zone and the model. Only
$\mathrm{adj}(v) = \emptyset$ needs $v$ expanded. So if either holds at an accepting $v$, then
$\Zn_v \subseteq \mathrm{stop}(v)$ and $W_v = \Zn_v$ is \emph{already} a post-fixpoint on
$S = \{v\}$: the query is settled with zero expansions and no fixpoint iteration.

They are therefore tested the instant a vertex is interned, not when it is eventually popped. On a
violated query whose witness is a single state, that is the difference between settling when the
state is \emph{found} and settling when the search gets round to it.

\subsection{Accepting states jump the queue, on a budget}
\label{sec:otf:order}

Only a $p \wedge \lnot q$ state can settle a violated query, so a newly generated one is the most
valuable thing in the frontier and need not wait behind the rest of its breadth level. The effect
is large where the first accepting state is already a witness: \texttt{phil4 P0.Eat --> P0.Think}
goes from $5$ outer expansions to $2$, because the process enumeration order puts the four fork
processes before the philosophers, so \texttt{P0}'s own move is generated last. \measured{}

Unbudgeted it is a trap. Where most accepting states are \emph{not} witnesses, front-loading them
turns the breadth search into a dive: \texttt{grc4 T0.near --> T0.in} goes from $182$ outer
expansions to $1762$. Prioritising only until two accepting states have been refuted keeps the
\texttt{phil4} gain (whose first accepting state never spends the budget) and restores the
\texttt{grc4} behaviour --- in fact improves it, to $118$. \measured{}

\subsection{Two exact restrictions}
\label{sec:otf:closure}

\textbf{The accepting closure.} $F$ is defined per vertex from its own successors, so the value at
an accepting vertex depends only on the $\lnot q$-forward-closure of the accepting set. Restricting
the final pass to that closure is \emph{exact}, not an approximation.

\textbf{The federation algebra.} $W_v \cap \mathrm{cand}$ was computed as a double complement,
$\mathrm{fed\_subtract}(z, \mathrm{fed\_subtract}(z, \mathrm{cand}))$. That denotes the right set,
but the inner subtraction peels $z$ by every constraint of every member of $\mathrm{cand}$ and
returns an \emph{arrangement} whose size is not bounded by $|\mathrm{cand}|$. The intersection of
two zones is a single zone, so the direct form needs exactly $|\mathrm{cand}|$ pieces. Compaction
additionally merges any pair whose union is exactly convex. Together these are the difference
between a federation that grows without bound and one that does not: \texttt{csma8 Bus.Active -->
Bus.Idle}, which did not converge at a cap of $2 \times 10^6$ in ten minutes, is settled in
$0.20$\,s with maximum federation width $1$. \measured{}
